%%%%%%%%%%%%%%%%%%%%%%%%%%%%%%%%%%%%%%%%%%%%%%%%%%%%%%%%%%%%%%%%%%%%%%%%%%%%%%%%
% ArTE-VLA draft for an intended ICRA 2027 submission.
% Internal drafting notes should remain in comments rather than visible prose.
% Before submission: add multi-seed results, complete ablations, replace the
% author block, add figures, and verify the final ICRA format and bibliography.
%%%%%%%%%%%%%%%%%%%%%%%%%%%%%%%%%%%%%%%%%%%%%%%%%%%%%%%%%%%%%%%%%%%%%%%%%%%%%%%%

\documentclass[letterpaper,10pt,conference]{ieeeconf}
\IEEEoverridecommandlockouts
\overrideIEEEmargins

\usepackage{amsmath}
\usepackage{amssymb}
\usepackage{booktabs}
\usepackage{graphicx}
\usepackage{url}
\usepackage{xcolor}

\newcommand{\method}{ArTE-VLA}
\newif\ifdraftnotes
\draftnotestrue % Set to \draftnotesfalse before submission.
\newcommand{\draftnote}[1]{%
  \ifdraftnotes
    \par\smallskip\noindent
    \parbox{\columnwidth}{%
      \color{red}\footnotesize\textbf{TODO:} #1}%
    \par\smallskip
  \fi
}

\title{\LARGE \bf
ArTE-VLA: Autoregressive Target Effects for Vision-Language-Action Models
}

\author{Anonymous Authors}

\begin{document}
\maketitle
\thispagestyle{empty}
\pagestyle{empty}

%%%%%%%%%%%%%%%%%%%%%%%%%%%%%%%%%%%%%%%%%%%%%%%%%%%%%%%%%%%%%%%%%%%%%%%%%%%%%%%%
\begin{abstract}
Vision-language-action (VLA) models typically generate robot actions directly
from visual observations and language instructions, leaving the intended
physical evolution of the scene implicit. We introduce \method{}, a VLA that
predicts how the language-referred target should change before generating an
action. Given dual-view RGB observations and an instruction, \method{} infers an
implicit target representation and autoregressively predicts target effects at
four future horizons. Each effect describes translation, rotation, linear and
angular velocity, and activity. The predicted effects are embedded as tokens and
appended to the visual-language prefix, allowing the action expert to condition
on a compact physical future. Simulator states provide effect supervision during
training but are not required at inference. In a single-seed study on all ten
LIBERO-Object tasks, \method{} improves the success rate of a recipe-matched
PI0.5 baseline from 70.8\% to 83.6\% over 500 paired episodes. Fixed-noise
interventions further show that changing the predicted effect values produces
consistent changes in generated actions. These results suggest that a compact,
target-centric future representation can provide an effective and inspectable
interface between VLA perception and control.
\end{abstract}

\draftnote{Recompute the abstract's numerical claims after multi-seed
replication. Until then, retain the explicit single-seed qualification.}

%%%%%%%%%%%%%%%%%%%%%%%%%%%%%%%%%%%%%%%%%%%%%%%%%%%%%%%%%%%%%%%%%%%%%%%%%%%%%%%%
\section{INTRODUCTION}
\label{sec:introduction}

Vision-language-action models have emerged as a promising route toward general
robot policies. By combining pretrained visual-language representations with
large collections of robot demonstrations, recent VLAs can follow open-vocabulary
instructions and transfer across objects, tasks, and environments
\cite{pi05}. Most existing architectures, however, are optimized to map the
current observation directly to an action sequence. The physical outcome that
connects an instruction to an action---which object should move, in what way,
and over what time scale---remains latent inside the policy.

This direct formulation creates two related difficulties. First, the action
expert must simultaneously resolve language grounding, infer task-relevant
dynamics, and generate precise controls. Second, the resulting computation is
difficult to inspect: a successful rollout does not reveal whether the policy
anticipated the intended object motion or relied on a less robust correlation.
Predictive world models and visual chain-of-thought methods address this problem
by forecasting future images or point trajectories before acting
\cite{cotvla,worldvla,atm,track2act}. Although expressive, these intermediate
representations can be high-dimensional and need not isolate the physical state
of the object that is central to the instruction.

We propose \method{}, which inserts a compact target-effect bottleneck into a
pretrained VLA. Rather than predicting a complete future image, the model
compresses its visual-language prefix into an implicit representation of the
primary target and predicts the target's physical effects at four horizons.
Each prediction contains a 3-D translation, a 3-D rotation, linear and angular
velocity, and an activity logit. The predictions form an autoregressive chain:
the effect at one horizon is used to predict the next. Horizon-specific hidden
states and effect values are then projected into tokens that condition the
flow-matching action expert.

The target-effect representation is learned with privileged supervision but
used without privileged inputs. We recover object states from the simulator
states stored in LIBERO demonstrations and construct aligned multi-horizon
effect labels. These labels appear only in the auxiliary prediction loss; the
action expert always receives model-predicted effects. At test time, \method{}
uses the same RGB observations, instruction, and robot state as vanilla PI0.5.
This separation allows structured physical supervision to shape the policy
without introducing an object detector, simulator state, task identifier, or
ground-truth target state into the deployed model.

Our main contributions are:
\begin{itemize}
    \item We introduce a compact, target-centric physical-effect bottleneck for
    VLA policies. The bottleneck predicts a four-horizon autoregressive future
    and directly conditions the action expert.
    \item We develop a replay-based labeling pipeline that adds aligned target
    effects to robot demonstrations while preserving their original images,
    actions, states, and language instructions.
    \item We formulate value-preserving interventions that separate reliance on
    the complete effect token from reliance on its explicit physical values and
    temporal order.
    \item On LIBERO-Object, \method{} improves a matched PI0.5 baseline by 12.8
    percentage points in a 500-episode paired evaluation, while revealing
    substantial variation across target objects.
\end{itemize}

%%%%%%%%%%%%%%%%%%%%%%%%%%%%%%%%%%%%%%%%%%%%%%%%%%%%%%%%%%%%%%%%%%%%%%%%%%%%%%%%
\section{RELATED WORK}
\label{sec:related_work}

\subsection{Vision-Language-Action Models}

VLA models adapt pretrained vision-language representations to continuous robot
control. PI0.5 combines heterogeneous co-training with a flow-matching action
expert to support open-world manipulation \cite{pi05}. \method{} uses PI0.5 as
its backbone and leaves the original visual-language prefix and action head
intact. Its contribution is an intermediate physical prediction interface that
is trained jointly with the action objective and consumed directly by the
action expert.

\subsection{Predictive Representations for Control}

World-model approaches learn scene dynamics through future observation
prediction. GR-2 jointly learns video generation and action prediction from web
videos and robot trajectories \cite{gr2}, while WorldVLA unifies image and
action generation within an autoregressive model \cite{worldvla}. CoT-VLA
generates future visual goals as an explicit visual chain of thought before
predicting actions \cite{cotvla}. These methods represent future scene content
with image or video tokens. In contrast, \method{} predicts a low-dimensional
physical state change for one language-conditioned target. The two approaches
are complementary: visual prediction preserves scene appearance, whereas a
target-effect bottleneck exposes a concise variable that can be measured and
intervened upon.

\subsection{Object and Point Motion as Policy Guidance}

Predicted motion can also serve as a control representation. Any-point
Trajectory Modeling predicts the future motion of arbitrary image points and
uses those trajectories to guide policy learning \cite{atm}. Track2Act predicts
goal-conditioned point tracks from internet videos and converts them into object
and end-effector motion plans \cite{track2act}. \method{} shares the premise that
future motion is useful for control, but differs in three respects: it predicts
a single instruction-grounded target, represents motion with multi-horizon 3-D
physical effects rather than 2-D tracks, and injects the prediction into the
token sequence of an end-to-end VLA action expert.

\draftnote{Expand and update the related-work review before submission,
including target-conditioned object representations, action-conditioned world
models, and contemporaneous VLA methods. Verify the final publication metadata
for every arXiv citation.}

%%%%%%%%%%%%%%%%%%%%%%%%%%%%%%%%%%%%%%%%%%%%%%%%%%%%%%%%%%%%%%%%%%%%%%%%%%%%%%%%
\section{METHODOLOGY}
\label{sec:method}

\subsection{Problem Formulation}

At time $t$, the policy receives RGB observations $I_t$ from a workspace and a
wrist camera, robot state $s_t$, and language instruction $\ell$. It predicts an
action chunk $A_t=(a_t,\ldots,a_{t+C-1})$. We augment this mapping with a target
effect chain $E_t=(e_t^1,\ldots,e_t^H)$ at $H=4$ future horizons:
\begin{align}
p(E_t\mid I_t,s_t,\ell)
&=\prod_{h=1}^{H}p(e_t^h\mid e_t^{<h},I_t,s_t,\ell), \\
p(A_t\mid I_t,s_t,\ell)
&=p_{\theta}(A_t\mid X_t,T(E_t)),
\label{eq:factorization}
\end{align}
where $X_t$ denotes the original PI0.5 prefix and $T(E_t)$ denotes the predicted
effect tokens. We use frame offsets $(1,4,8,16)$, corresponding to $(0.1,0.4,
0.8,1.6)$ seconds at the 10-Hz control rate. Autoregression in
Eq.~\eqref{eq:factorization} refers to temporal prediction and does not imply
identification of a structural causal model.

\subsection{Implicit Target Representation}

PI0.5 encodes the two RGB observations and instruction into prefix embeddings
$X_t\in\mathbb{R}^{N\times d}$. We apply masked mean pooling over the prefix and
project the result to a 20-D vector:
\begin{equation}
z_t=f_{\mathrm{target}}\!\left(
\frac{\sum_{i=1}^{N}m_iX_{t,i}}{\sum_{i=1}^{N}m_i}
\right)\in\mathbb{R}^{20}.
\label{eq:target}
\end{equation}
The vector $z_t$ is the sole active entity input to the effect predictor. It is
not directly regressed to the current simulator state; instead, its semantics
are induced by language-conditioned future-effect supervision and the action
loss. We therefore treat it as an implicit target representation rather than a
calibrated estimate of object pose.

\subsection{Multi-Horizon Effect Chain}

The effect model projects $z_t$ into a hidden state and combines it with a
pooled VLA context, a learned horizon embedding, and the previously predicted
effect. At horizon $h$, it produces
\begin{equation}
e_t^h=[\Delta p_t^h,\Delta r_t^h,v_t^h,\omega_t^h,q_t^h]
\in\mathbb{R}^{13},
\label{eq:effect}
\end{equation}
where $\Delta p_t^h\in\mathbb{R}^{3}$ and
$\Delta r_t^h\in\mathbb{R}^{3}$ describe translation and axis-angle rotation
relative to time $t$, $v_t^h,\omega_t^h\in\mathbb{R}^{3}$ are future linear and
angular velocities, and $q_t^h$ is an activity logit. The prediction $e_t^h$ is
projected back into the hidden space when computing $e_t^{h+1}$.

\subsection{Conditioning the Action Expert}

The effect predictor retains a horizon-specific hidden context $c_t^h$. We
combine this context with the corresponding physical prediction and project the
result to the PI0.5 embedding dimension:
\begin{equation}
\tau_t^h=f_{\mathrm{token}}([c_t^h;e_t^h]).
\label{eq:token}
\end{equation}
The four effect tokens are appended to the original prefix,
$\widetilde{X}_t=[X_t;\tau_t^1;\ldots;\tau_t^4]$, and are visible to the
flow-matching action expert. The original prefix is never replaced. Consequently,
the policy can use the predicted physical future while retaining the full visual
and linguistic context of the backbone.

\subsection{Training Objective}

For each valid horizon, simulator replay supplies a target effect
$\hat e_t^h$ and mask $M_t^h$. The loss is
\begin{equation}
\mathcal{L}=\mathcal{L}_{\mathrm{action}}
+\lambda_E\bigl(\mathcal{L}_{\mathrm{reg}}+\mathcal{L}_{\mathrm{act}}\bigr),
\label{eq:objective}
\end{equation}
where $\mathcal{L}_{\mathrm{reg}}$ is masked Smooth-$L_1$ loss over the 12
continuous values, $\mathcal{L}_{\mathrm{act}}$ is binary cross-entropy over
activity, and $\lambda_E=1$. We train the effect chain without teacher forcing:
the action expert is conditioned on model predictions, never on future labels.

\draftnote{Expand Eq.~\eqref{eq:objective} to specify horizon aggregation, mask
normalization, active/inactive weighting, and the exact action-loss definition
implemented in the training code.}

\subsection{Replay-Based Effect Supervision}

We construct effect labels by restoring the flattened simulator state stored at
each demonstration frame and forwarding the simulator without re-executing the
recorded action. For the supervised target, we measure pose displacement and
future velocity at each offset in Eq.~\eqref{eq:factorization}. Horizons that
cross an episode boundary are masked. The conversion preserves the original RGB
observations, robot state, action, instruction, task index, episode index, and
frame index. On LIBERO-Object, the target occupies the same entity slot across
all ten tasks. The task index is used only to select a training label and is not
embedded by the policy or provided during rollout.

\subsection{Effect Interventions}

We intervene on predicted effects while fixing the observation, language,
checkpoint, diffusion noise, and denoising schedule. \emph{Mask} removes all
effect tokens. \emph{Zero-12D} sets the continuous values in
Eq.~\eqref{eq:effect} to zero while preserving the horizon context and activity
logit. \emph{Reverse} reverses the order of the four effect values, and
\emph{shuffle} substitutes values from another sample while retaining the
receiving sample's hidden contexts. These interventions distinguish dependence
on the added token pathway from dependence on predicted physical values and
their temporal organization.

\draftnote{Add the main architecture figure showing the PI0.5 backbone, implicit
target representation, four autoregressive effect steps, action conditioning,
and the training-only simulator-supervision pathway.}

%%%%%%%%%%%%%%%%%%%%%%%%%%%%%%%%%%%%%%%%%%%%%%%%%%%%%%%%%%%%%%%%%%%%%%%%%%%%%%%%
\section{EXPERIMENTS}
\label{sec:experiments}

\subsection{Benchmark and Data}

We evaluate on LIBERO-Object, which contains ten language-conditioned tabletop
manipulation tasks \cite{libero}. The training set contains 454 available
demonstration episodes and 66,984 frames. Replay enrichment yields 253,423 valid
target-horizon labels, of which 170,107 are active. We verify frame-level
alignment between the original and enriched datasets and observe a maximum
action-alignment error of $2.90\times10^{-8}$. Terminal horizons without a valid
successor state are masked.

\subsection{Training Setup}

Both vanilla PI0.5 and \method{} are initialized from
\texttt{lerobot/pi05\_base} and fine-tuned for 30,000 updates using eight NVIDIA
H20 GPUs, global batch size 64, bfloat16 precision, and seed 1000. We use the same
optimizer, learning-rate schedule, image resolution, action chunk length of 16,
and five executed actions per inference call. \method{} adds the target head,
effect chain, effect-token projection, and the auxiliary loss in
Eq.~\eqref{eq:objective}; the remaining training recipe is unchanged.

\draftnote{Add the exact optimizer, learning rate and schedule, weight decay,
image resolution, augmentation, checkpoint-selection rule, wall-clock training
cost, and trainable-parameter overhead for reproducibility.}

\subsection{Evaluation Protocol}

For each task, we evaluate both policies from the same 50 demonstration-derived
initial simulator states, resulting in 500 paired episodes per method. Rollouts
use synchronous batches of ten environments, ten flow-matching inference steps,
and a fixed time horizon. Success is recorded without terminating the episode
early. We report per-task and aggregate success, paired rescue and regression
counts, and the exact two-sided McNemar test. Because these resets reconstruct
expert demonstration states rather than the benchmark's standard initial-state
set, we keep this protocol separate from comparisons to previously published
LIBERO results.

\subsection{Baselines and Ablations}

The primary baseline is recipe-matched vanilla PI0.5. To distinguish future
prediction from auxiliary regularization and added capacity, the complete study
also considers an effect predictor that does not condition actions, learned
context tokens without structured values, current-target-state conditioning,
non-autoregressive horizon prediction, and a privileged ground-truth-effect
oracle. We evaluate removal, zero-value, reverse-order, and cross-sample shuffle
interventions on the same \method{} checkpoint. These comparisons isolate three
questions: whether the additional pathway improves control, whether future
structure matters beyond additional tokens, and whether the trained action
expert uses the predicted values.

\draftnote{Run and report every baseline and ablation named above, including
capacity-matched and value-free controls. Remove any comparison from the prose
if it is not completed.}

\subsection{Metrics}

Closed-loop success is the primary metric. For effect prediction, we measure
translation average and final displacement error, rotation error, linear and
angular velocity error, and activity precision and recall. We also
report prediction improvement over a zero-motion baseline. Intervention effects
are measured by the fraction of changed action chunks, maximum absolute action
difference, and relative root-mean-square (RMS) distance. Efficiency is measured
with trainable parameter count, inference latency, and peak accelerator memory.

\draftnote{Add the effect-prediction and efficiency results promised by this
section, with per-horizon errors, a zero-motion comparison, parameter overhead,
latency, and peak memory; otherwise remove unreported metrics.}

%%%%%%%%%%%%%%%%%%%%%%%%%%%%%%%%%%%%%%%%%%%%%%%%%%%%%%%%%%%%%%%%%%%%%%%%%%%%%%%%
\section{RESULTS}
\label{sec:results}

\subsection{Closed-Loop Performance}

Table~\ref{tab:main_results} compares \method{} with vanilla PI0.5 and reserves
matched columns for the auxiliary-only and value-free-token baselines.
Aggregated over all ten tasks, target-effect conditioning improves success from
70.8\% to 83.6\%, an absolute gain of 12.8 percentage points. Across the 500
paired episodes, \method{} produces 127 rescues and 63 regressions. The
corresponding exact two-sided McNemar test yields $p=3.99\times10^{-6}$.

\begin{table*}[t]
\caption{LIBERO-Object success over 50 paired episodes per task. Available
results are from one training seed. Blank cells indicate matched baselines
awaiting training and evaluation.}
\label{tab:main_results}
\centering
\small
\setlength{\tabcolsep}{9pt}
\begin{tabular}{lrrrrr}
\toprule
Target object & PI0.5 & Auxiliary-only & Value-Free Tokens & \method{} & $\Delta$ vs. PI0.5 \\
\midrule
Alphabet Soup       & 76\% & & & 74\% & $-2$ \\
Cream Cheese        & 76\% & & & 94\% & $+18$ \\
Salad Dressing      & 88\% & & & 82\% & $-6$ \\
BBQ Sauce           & 96\% & & & 84\% & $-12$ \\
Ketchup             & 98\% & & & 64\% & $-34$ \\
Tomato Sauce        & 74\% & & & 92\% & $+18$ \\
Butter              &  4\% & & & 94\% & $+90$ \\
Milk                & 96\% & & & 70\% & $-26$ \\
Chocolate Pudding   &  6\% & & & 84\% & $+78$ \\
Orange Juice        & 94\% & & & 98\% & $+4$ \\
\midrule
\textbf{Overall}    & \textbf{70.8\%} & & & \textbf{83.6\%} & \textbf{$+12.8$} \\
\bottomrule
\end{tabular}
\end{table*}

\draftnote{Repeat the main comparison with additional training seeds and report
uncertainty intervals. Add a separate standard-LIBERO-initialization evaluation;
do not merge it with the expert-demonstration-reset protocol used in Table~I.}

\draftnote{Add a compact results visualization or paired rescue/regression
summary so that the strongly heterogeneous per-task behavior is visible without
relying only on aggregate success.}

The aggregate improvement is dominated by two large rescues. On Butter and
Chocolate Pudding, vanilla PI0.5 succeeds in only 4\% and 6\% of episodes,
whereas \method{} reaches 94\% and 84\%, respectively. The gain is therefore
not a uniform shift across tasks. Ketchup and Milk decrease by 34 and 26 points,
and BBQ Sauce decreases by 12 points. This heterogeneity indicates that target
effects can resolve severe failures of a direct policy, but can also introduce
task-dependent interference.

\subsection{Component Ablations}

Table~\ref{tab:ablations} isolates the contribution of action conditioning,
structured effect values, present-versus-future target information, and
autoregressive horizon modeling. Each learned variant will use the same
initialization, training budget, and evaluation protocol as the full model.

\begin{table}[t]
\caption{Matched component ablations on LIBERO-Object. Blank cells denote
experiments that have not yet been completed.}
\label{tab:ablations}
\centering
\small
\begin{tabular}{lrr}
\toprule
Variant & Success & $\Delta$ vs. PI0.5 \\
\midrule
Auxiliary-only              & & \\
Value-Free Tokens           & & \\
Current-target conditioning & & \\
Non-autoregressive effects  & & \\
\midrule
Full \method{}              & 83.6\% & $+12.8$ \\
\bottomrule
\end{tabular}
\end{table}

\draftnote{Populate Table~\ref{tab:ablations} with separately trained,
recipe-matched checkpoints and multi-seed uncertainty intervals.}

\subsection{Action Dependence on Effect Values}

We perform a fixed-noise action audit on 32 observations drawn equally from
Butter and Chocolate Pudding. Repeating the unmodified forward computation gives
zero numerical action difference. Zeroing the 12 continuous effect values
changes the action chunk for all 32 samples. The median maximum absolute action
difference is 0.00956, and the median relative RMS distance is 0.943\%.
Reverse-order and cross-sample shuffle interventions also produce consistent
action changes under the fixed-noise protocol.

\begin{table}[t]
\caption{Fixed-noise effect interventions on 32 observations. Action metrics
are measured relative to the unmodified computation. Blank cells are pending;
closed-loop success requires full-task evaluation.}
\label{tab:interventions}
\centering
\footnotesize
\setlength{\tabcolsep}{3.5pt}
\begin{tabular}{lrrrr}
\toprule
Setting & Changed & Max abs. & Rel. RMS & Success \\
        & (\%)    &          & (\%)    & (\%) \\
\midrule
Unmodified repeat & 0   & 0       & 0     & \\
Mask              &     &         &       & \\
Zero-12D          & 100 & 0.00956 & 0.943 & \\
Reverse           &     &         &       & \\
Shuffle           &     &         &       & \\
\bottomrule
\end{tabular}
\end{table}

These results establish computational dependence of the action expert on the
predicted physical values. They do not establish that those values alone account
for the closed-loop improvement: each effect token also contains a learned
horizon context, and the zero-value median lies slightly below the prespecified
1\% relative-RMS decision threshold. Closed-loop interventions over the full
task distribution are therefore necessary to quantify the independent control
value of the structured prediction.

\draftnote{Extend the value intervention beyond the current 32 observations and
two tasks. Run all-task action audits and closed-loop mask, zero-value,
reverse-order, and cross-sample-shuffle interventions before making an
independent value-use claim.}

\subsection{Task-Level Behavior}

The paired outcomes suggest two complementary failure modes. For tasks on which
vanilla PI0.5 rarely succeeds, target-effect supervision may provide an
additional object-motion signal that disambiguates action generation. Conversely,
the regressions on Ketchup, Milk, and BBQ Sauce may arise from inaccurate target
effects, an implicit target representation that entangles object and task
context, or competition between effect and visual-language tokens inside the
action expert. Separating these explanations requires comparing effect error and
action sensitivity on matched rescue and regression episodes.

\draftnote{Add the matched rescue/regression diagnostic proposed above,
including effect-error correlations and representative qualitative rollouts;
otherwise keep these failure explanations explicitly speculative.}

%%%%%%%%%%%%%%%%%%%%%%%%%%%%%%%%%%%%%%%%%%%%%%%%%%%%%%%%%%%%%%%%%%%%%%%%%%%%%%%%
\section{DISCUSSION}
\label{sec:discussion}

The current results support a narrow conclusion: adding an autoregressive
target-effect pathway to PI0.5 improves aggregate success on LIBERO-Object and
measurably changes the action computation. They do not yet show that predicting
future physical values is uniformly beneficial across tasks or that explicit
values are the sole source of improvement. In particular, learned horizon
contexts provide an additional conditioning channel that must be controlled with
parameter-matched and value-free baselines.

The method also relies on privileged target-effect labels during training. This
supervision is compatible with simulation and instrumented robot datasets, but
application to unconstrained real-world data would require state estimation or
pseudo-labeling. In addition, the current target representation is implicit and
does not expose a calibrated present-time object pose. An explicit target
estimator could improve interpretability, but would introduce association and
geometry errors that the present architecture avoids.

Finally, Table~\ref{tab:main_results} reports a single training seed and an
expert-demonstration reset protocol. Multi-seed replication, matched auxiliary
and capacity controls, standard LIBERO initial states, and evaluation beyond
LIBERO-Object are required to determine whether the observed gain is robust and
generalizes beyond the current setting.

\draftnote{Add standard LIBERO initial-state results, held-out task suites, and
cross-simulator evaluation, or narrow the final generalization claims to the
current LIBERO-Object expert-reset setting.}

%%%%%%%%%%%%%%%%%%%%%%%%%%%%%%%%%%%%%%%%%%%%%%%%%%%%%%%%%%%%%%%%%%%%%%%%%%%%%%%%
\section{CONCLUSION}
\label{sec:conclusion}

We presented \method{}, a VLA architecture that predicts a compact physical
future for the language-referred target and uses that prediction to condition
action generation. The method requires privileged object states only for
training-time supervision and retains the original RGB-language interface at
inference. On LIBERO-Object, \method{} substantially improves aggregate success
over a matched PI0.5 baseline, while intervention experiments confirm that the
action expert responds to the predicted values. The task-level regressions and
the contribution of hidden effect context identify important limitations. More
broadly, our results motivate structured future effects as an inspectable
interface between visual-language understanding and continuous robot control.

\draftnote{Revise the conclusion after the outstanding replications, ablations,
and interventions; ensure that no final claim exceeds the completed evidence.}

\draftnote{Complete the bibliography audit: expand coverage, replace arXiv
entries with archival publication data where available, and verify author lists,
titles, years, and identifiers.}

%%%%%%%%%%%%%%%%%%%%%%%%%%%%%%%%%%%%%%%%%%%%%%%%%%%%%%%%%%%%%%%%%%%%%%%%%%%%%%%%
\begin{thebibliography}{99}

\bibitem{pi05}
Physical Intelligence et al., ``$\pi_{0.5}$: A vision-language-action model with
open-world generalization,'' arXiv preprint arXiv:2504.16054, 2025.

\bibitem{libero}
B. Liu, Y. Zhu, C. Gao, Y. Feng, Q. Liu, Y. Zhu, and P. Stone, ``LIBERO:
Benchmarking knowledge transfer for lifelong robot learning,'' in
\emph{Advances in Neural Information Processing Systems}, vol. 36, 2023.

\bibitem{cotvla}
Q. Zhao et al., ``CoT-VLA: Visual chain-of-thought reasoning for
vision-language-action models,'' arXiv preprint arXiv:2503.22020, 2025.

\bibitem{worldvla}
J. Cen et al., ``WorldVLA: Towards autoregressive action world model,'' arXiv
preprint arXiv:2506.21539, 2025.

\bibitem{gr2}
C.-L. Cheang et al., ``GR-2: A generative video-language-action model with
web-scale knowledge for robot manipulation,'' arXiv preprint arXiv:2410.06158,
2024.

\bibitem{atm}
C. Wen, X. Lin, J. So, K. Chen, Q. Dou, Y. Gao, and P. Abbeel, ``Any-point
trajectory modeling for policy learning,'' arXiv preprint arXiv:2401.00025,
2024.

\bibitem{track2act}
H. Bharadhwaj, R. Mottaghi, A. Gupta, and S. Tulsiani, ``Track2Act: Predicting
point tracks from internet videos enables generalizable robot manipulation,''
arXiv preprint arXiv:2405.01527, 2024.

\end{thebibliography}

\draftnote{Before submission, set \texttt{\string\draftnotesfalse}, confirm the
anonymous author block for review, remove identifying metadata and comments,
and rerun the PaperPlaza PDF compliance check.}

\end{document}
